% sashimi 2026 -- ext a (patchnce hybrid), 8pp + up to 2pp refs LNCS, double-blind.
% VERIFY the exact 2026 page limit on sashimiworkshop.org/submission.
% anonymized: no names, emails, institutions, grant ids, repo urls.
% framing: a MECHANISM STUDY (complementary failure modes of cycle vs
% patch-contrastive content preservation) with corrected metrics + downstream
% validation -- NOT a "we added a loss" paper. numerics marked TODO pending the
% corrected evaluation. do NOT headline the global-moment ssim.

\documentclass[runningheads]{llncs}

\usepackage{graphicx}
\usepackage{amsmath,amssymb,bm}
\usepackage{booktabs}
\usepackage{multirow}
\usepackage{subcaption}
\usepackage[hidelinks]{hyperref}
\usepackage{xcolor}
\newcommand{\todo}[1]{\textcolor{red}{[TODO: #1]}}

\begin{document}

\title{Cycle-Consistent Patch-Contrastive Harmonization of\\
Multi-Site Glioma MRI: A Mechanism Study\\with Downstream Validation}

\titlerunning{Cycle-Consistent Patch-Contrastive MRI Harmonization}

\author{Anonymous}
\authorrunning{Anonymous (SASHIMI 2026 submission)}
\institute{Paper ID \#TODO. Author and institution details omitted for double-blind review.}

\maketitle

\begin{abstract}
Scanner- and protocol-induced intensity shift confounds multi-institutional MRI
studies, and unpaired image translation is an increasingly common remedy.
Two largely disjoint mechanisms preserve anatomical content during translation:
\emph{global, bidirectional} cycle-consistency, and \emph{local, feature-space}
patch-contrastive alignment (PatchNCE). We argue these mechanisms fail in
complementary ways under unpaired glioma data---the cycle constraint admits an
appearance-leaking, boundary-smoothing equilibrium, while one-sided PatchNCE
lacks a global appearance anchor---and we study their combination on a
self-attention 2.5D tri-planar backbone. We frame the work as a \emph{mechanism
study}: rather than asking ``does adding PatchNCE help,'' we ask under what
weighting and measured by which task the cycle/contrastive combination resolves
a failure mode that neither component resolves alone. On a
BraTS\,$\leftrightarrow$\,UPenn-GBM benchmark we evaluate with distribution
metrics (FID/KID, domain-classifier confusion), brain-masked \emph{windowed}
structural similarity, and \emph{downstream} cross-site glioma sub-region
segmentation (Dice/HD95), ablating the contrastive weight
$\lambda_{\mathrm{NCE}}\in\{0,0.1,0.5,1.0,2.0\}$ at a matched fine-tuning budget.
\todo{headline result sentence once corrected metrics land.}
\keywords{MRI harmonization \and unpaired translation \and contrastive learning \and cycle-consistency \and downstream validation.}
\end{abstract}

\section{Introduction}
\label{sec:intro}
Pooling magnetic resonance imaging (MRI) across institutions is essential for
statistically powered neuro-oncology studies, yet scanner vendor, field strength,
receive-coil configuration, and acquisition protocol induce systematic intensity
and contrast variation that confounds quantitative analysis. Feature-space
correction such as ComBat~\cite{fortin2017combat} removes batch effects from
\emph{extracted features} but cannot produce harmonized \emph{images}, and is
therefore unavailable to spatial downstream tasks---tumor segmentation, lesion
volumetry, voxel-level radiomics---that operate on the image itself. This
motivates generative, image-space harmonization, in which an unpaired translation
model maps an image acquired at one site to the appearance of another while
preserving anatomy.

The central difficulty is preserving content. Two mechanisms dominate the
unpaired-translation literature. \emph{Cycle-consistency}~\cite{zhu2017cyclegan}
enforces content preservation \emph{globally and bidirectionally} through
pixel-space reconstruction, but the cycle constraint is satisfiable by an
appearance-leaking equilibrium: under unpaired data with pathology, the generator
can hide source information in low-amplitude, high-frequency signal that survives
the round trip, manifesting as over-smoothed tissue boundaries and, in glioma
cohorts, attenuated or hallucinated lesion texture. \emph{Patch-contrastive}
translation (CUT/PatchNCE~\cite{park2020cut}) instead preserves content
\emph{locally in feature space} by maximizing patch-level mutual information,
which is robust to the cycle's failure mode---but applied one-sided it provides
no global anchor on the target appearance distribution and is known to require
aggressive weighting and careful negative sampling to remain
stable~\cite{wang2021negcut,zhan2022monce}. We hypothesize that these failure
modes are \emph{complementary}, and we study their combination on a
self-attention 2.5D tri-planar backbone for within-modality glioma harmonization.

A second, methodological point motivates this work. Harmonization quality is
routinely reported with image-similarity scores, but a global structural-similarity
index over an MRI slice is dominated by the large background region and is
misleadingly high; worse, it does not reveal whether disease-relevant signal is
preserved. The appropriate test is \emph{downstream}: does a model trained at one
site retain its accuracy when applied, after harmonization, to another? We
therefore center our evaluation on cross-site tumor-segmentation transfer, and we
report distribution-matching and \emph{brain-masked windowed} similarity rather
than a single global score.

\paragraph{Contributions.}
\begin{enumerate}
  \item A \emph{mechanism} analysis of cycle-consistency versus patch-contrastive
        content preservation for within-modality MRI harmonization, with a
        controlled $\lambda_{\mathrm{NCE}}$ ablation (including the
        $\lambda_{\mathrm{NCE}}{=}0$ cycle-only control) at a matched fine-tuning
        budget that locates the operating regime in which the combination is
        beneficial.
  \item A downstream-validated, training-consistent evaluation protocol on glioma
        cohorts (BraTS/TCGA-GBM, UPenn-GBM): cross-site sub-region segmentation
        Dice/HD95, distribution-matching (FID/KID, domain-classifier confusion),
        and brain-masked windowed similarity. We show that a naive
        harmonization-for-evaluation pipeline (mismatched channel order and
        normalization) silently corrupts results, and we describe the consistent
        protocol that avoids it.
  \item \todo{the empirical finding, stated honestly once corrected metrics land
        (e.g.\ the regime in which the hybrid improves cross-site Dice, and the
        correction of a previously-reported downstream regression).}
\end{enumerate}

\section{Related Work}
\label{sec:related}
\paragraph{Contrastive unpaired translation.}
CUT~\cite{park2020cut} replaced cycle-consistency with a one-sided multi-layer
PatchNCE objective drawing negatives from within the input image;
DCLGAN~\cite{han2021dclgan} showed contrastive and cycle-style dual objectives
are complementary on natural images. Subsequent work refined where and how
patches are sampled---query selection~\cite{hu2022qsattn}, semantic-relation
consistency~\cite{jung2022src}, hard-negative generation~\cite{wang2021negcut},
and modulated negatives~\cite{zhan2022monce}---all instances of the InfoNCE
estimator~\cite{oord2018cpc}. We deliberately \emph{re-introduce} full
cycle-consistency on top of multi-layer PatchNCE---the inverse of CUT's design
choice---and justify it rather than hide it: cycle supplies the global appearance
anchor that one-sided PatchNCE lacks, so the contrastive term can act as a
complement at a weight well below the CUT/FastCUT regime.

\paragraph{MRI harmonization.}
Beyond feature-space ComBat~\cite{fortin2017combat}, image-space harmonizers
include disentanglement methods (CALAMITI~\cite{zuo2021calamiti},
HACA3~\cite{zuo2023haca3}), the 2.5D VAE-GAN ImUnity~\cite{cackowski2023imunity},
and the many-to-one CycleGAN IGUANe~\cite{roca2024iguane}. The closest
methodological precedent, cycleCUT~\cite{wang2022cyclecut}, combines cycle and
contrastive losses---but for cross-modality MR$\to$CT synthesis, with no
within-modality harmonization and no lesion-level downstream task; one-sided
contrastive medical translation otherwise relies on auxiliary segmentation
supervision~\cite{seibel2024oct}. To our knowledge no prior work studies a
cycle\,$+$\,multi-layer-PatchNCE hybrid for within-modality MRI harmonization
with glioma downstream validation. We position our work as a mechanism study in
that gap rather than as a new loss.

\paragraph{Evaluation.}
Harmonization is commonly assessed by image similarity and distribution distance.
We follow the stronger conventions: distribution matching via FID and
KID~\cite{heusel2017fid,binkowski2018kid}, a domain classifier driven toward
chance, brain-masked \emph{windowed} similarity rather than a global moment, and
downstream segmentation with a frozen cross-site segmenter
(Dice/HD95)~\cite{isensee2021nnunet}, all aggregated at the subject level.

\section{Method}
\label{sec:method}
\paragraph{Backbone.}
We build on a self-attention CycleGAN with a 2.5D tri-planar generator: each
input stacks three consecutive axial slices $(s_{z-1},s_z,s_{z+1})$ across four
modalities (FLAIR, T1, T1ce/T1gd, T2), giving inter-slice context at the memory
cost of a 2D model, with self-attention to capture the spatially non-stationary
component of scanner bias. The generators $G_{A\to B},G_{B\to A}$ output the
harmonized center slice; the architecture is summarized in
Fig.~\ref{fig:method}.

\begin{figure}[t]
  \centering
  \fbox{\parbox{0.96\linewidth}{\centering\vspace{2.2cm}\todo{method/architecture
  diagram: backbone + per-layer PatchNCE heads; cycle path (global appearance
  anchor) vs contrastive path (local feature-space content). reviewers of the
  prior submission explicitly flagged the absence of a method diagram.}\vspace{2.2cm}}}
  \caption{Self-attention 2.5D CycleGAN with per-layer PatchNCE projection heads.
  Cycle-consistency provides a global, bidirectional appearance/anatomy anchor;
  multi-layer PatchNCE provides a local, feature-space content regularizer.}
  \label{fig:method}
\end{figure}

\paragraph{Cycle-consistency (global anchor).}
For unpaired collections $\mathcal{D}_A,\mathcal{D}_B$,
\begin{equation}
  \mathcal{L}_{\mathrm{cyc}} =
  \mathbb{E}_{x_a}\!\big[\|x_a-G_{B\to A}(G_{A\to B}(x_a))\|_1\big]
  +\mathbb{E}_{x_b}\!\big[\|x_b-G_{A\to B}(G_{B\to A}(x_b))\|_1\big].
  \label{eq:cyc}
\end{equation}

\paragraph{Multi-layer PatchNCE (local anchor).}
Encoder features at $L$ layers are passed through 2-layer MLP heads $h_\ell$ and
L2-normalized. For a query $z$, its positive $z^{+}$ at the corresponding spatial
location, and internal negatives $z^{-}$, the InfoNCE loss is
\begin{equation}
  \ell(z,z^{+},z^{-}) = -\log
  \frac{\exp(z\!\cdot\!z^{+}/\tau)}
       {\exp(z\!\cdot\!z^{+}/\tau)+\sum_{n}\exp(z\!\cdot\!z^{-}_{n}/\tau)},
  \quad
  \mathcal{L}_{\mathrm{NCE}}=\mathbb{E}\sum_{\ell=1}^{L}\sum_{s}\ell(\cdot),
  \label{eq:nce}
\end{equation}
with temperature $\tau{=}0.07$, $L{=}5$ encoder stages, and $256$ sampled
locations per layer, following CUT~\cite{park2020cut}.

\paragraph{Hybrid objective.}
PatchNCE \emph{augments}, rather than replaces, cycle-consistency:
\begin{equation}
  \mathcal{L} = \mathcal{L}_{\mathrm{adv}} + \mathcal{L}_{\mathrm{cyc}}
  + \mathcal{L}_{\mathrm{idt}} + \mathcal{L}_{\mathrm{ssim}}
  + \lambda_{\mathrm{NCE}}\,\mathcal{L}_{\mathrm{NCE}}.
  \label{eq:total}
\end{equation}
We sweep $\lambda_{\mathrm{NCE}}\in\{0,0.1,0.5,1.0,2.0\}$ at otherwise identical
hyperparameters, where $\lambda_{\mathrm{NCE}}{=}0$ is the cycle-only control.
A selected weight below the CUT/FastCUT regime ($1.0$/$10.0$) is principled
rather than incidental: because cycle-consistency already supplies a strong
content constraint, the contrastive term should complement it, not dominate;
over-weighting re-introduces the instability that PatchNCE exhibits when it must
carry content alone.

\section{Experimental Setup}
\label{sec:exp}
\paragraph{Data.}
Site~$A$ is BraTS/TCGA-GBM ($88$ subjects) and site~$B$ is UPenn-GBM ($566$
subjects), each with four co-registered modalities and expert tumor
segmentations at $1\,\mathrm{mm}$ isotropic, intensity-normalized to $[0,1]$. We
split \emph{by subject} $80/10/10$ and sample axial 2.5D triplets, yielding
\todo{$42{,}110$ train / $5{,}263$ val / $5{,}265$ test} slices; all metrics are
aggregated to the subject level.
\paragraph{Training.}
Each $\lambda_{\mathrm{NCE}}$ arm is warm-started from a converged base
SA-CycleGAN and fine-tuned for a matched budget (Adam, $\beta_1{=}0.5$; mixed
precision; A100), so that any difference across arms reflects the contrastive
weight and not the optimization budget.
\paragraph{Metrics.}
Distribution matching: FID and KID (harmonized vs target, per modality), domain
classifier accuracy toward chance, and feature-space MMD. Structure: brain-masked
\emph{windowed} SSIM (we report windowed, masked values and explain why a global
moment is misleading for harmonization). Downstream: a frozen U-Net trained on one
site and applied cross-site, raw versus harmonized, reporting whole-tumor and
sub-region Dice and HD95. \emph{Critically, the harmonization used for downstream
evaluation reuses the exact training data pipeline} (channel order, $[0,1]$
normalization, slice-major 2.5D interleaving); we found that a mismatched
evaluation pipeline corrupts the generator input and can manufacture a spurious
downstream regression.
\paragraph{Baselines / ablation.}
The $\lambda_{\mathrm{NCE}}$ sweep includes the cycle-only control
($\lambda_{\mathrm{NCE}}{=}0$) on the same backbone, data, and budget; we compare
against the base SA-CycleGAN and \todo{CUT-only if compute permits}, with
identical training budget across arms.

\section{Results}
\label{sec:results}
\todo{ALL NUMBERS PENDING the corrected evaluation. Tables: (1) distribution +
windowed-masked structure metrics per $\lambda$ and vs cycle-only; (2) downstream
cross-site Dice/HD95 raw vs harmonized, both directions. Figures: $\lambda$ vs
downstream Dice; qualitative source/harmonized/target with a boundary/edge metric
showing the over-smoothing failure mode; failure cases. Report mean$\pm$std with
paired Wilcoxon across subjects. Regenerate from logged run dirs; never hand-edit.}

\section{Discussion and Limitations}
\label{sec:discuss}
\todo{honest reading of the mechanism question and the complementarity evidence;
correction of the prior downstream regression; limitations: single seed per arm,
two-site scope, slice budget, brain-mask source, CUT-only baseline if omitted.}

\section{Conclusion}
\label{sec:conclusion}
\todo{one short paragraph.}

\bibliographystyle{splncs04}
\bibliography{refs}

\end{document}
